\documentclass[12pt,letterpaper]{article}
\usepackage[utf8]{inputenc}
\usepackage[T1]{fontenc}
\usepackage{lmodern,amsmath,amssymb,array,enumitem,xcolor,geometry,tcolorbox}
\geometry{left=1.7cm,right=1.7cm,top=1.5cm,bottom=1.6cm}
\setlength{\parindent}{0pt}
\definecolor{azul}{HTML}{17365D}
\definecolor{claro}{HTML}{EAF1F8}
\newcommand{\familia}[2]{\section*{Familia #1. #2}}
\newcommand{\campo}[1]{\textcolor{azul}{\textbf{#1}}}
\begin{document}
\begin{center}
{\LARGE\bfseries\color{azul} SUCESIONES Y LÍMITES}\par
{\Large Banco paramétrico -- Nivel experto}
\end{center}
\begin{tcolorbox}[colback=claro,colframe=azul]
El nivel experto conserva los mismos aprendizajes, pero exige elegir una estrategia,
combinar técnicas o demostrar una propiedad no evidente. Los datos siguen siendo
pequeños; la dificultad proviene del razonamiento y no de cálculos tediosos.
\end{tcolorbox}

\familia{1}{Reconstrucción desde diferencias}
\campo{Plantilla:}
\[
 a_{n+1}-a_n=pn+q,\qquad a_1=r.
\]
Pedir hallar una fórmula explícita y verificarla por inducción.\par
\campo{Condiciones:} $p\neq0$, $|p|\leq4$, $|q|\leq5$, $|r|\leq6$.
Usar
\[
 a_n=r+\frac{p(n-1)n}{2}+q(n-1).
\]
\campo{Tuplas $(p,q,r)$:}
\[
 (2,1,3),\ (3,-1,2),\ (-2,5,1),\ (4,-3,-1).
\]

\familia{2}{Recurrencia de segundo orden}
\campo{Plantilla:}
\[
 b_n=A\alpha^n+B\beta^n.
\]
Pedir construir la recurrencia, calcular datos iniciales y demostrar equivalencia.\par
\campo{Condiciones:} $\alpha,\beta\in\{-2,-1,2,3,4\}$, $\alpha\neq\beta$,
$A,B\in\{-3,-2,-1,1,2,3\}$ y términos iniciales de valor absoluto menor que $50$.\par
\campo{Control:}
\[
 b_{n+2}=(\alpha+\beta)b_{n+1}-\alpha\beta b_n.
\]
\campo{Tuplas $(A,B,\alpha,\beta)$:}
\[
 (1,-1,3,2),\ (2,-1,2,-1),\ (1,2,2,-2),\ (-1,2,3,1).
\]

\familia{3}{Monotonía con cambio de comportamiento}
\campo{Plantilla:}
\[
 c_n=\frac{p^n}{(n+q)!}.
\]
Pedir localizar exactamente dónde crece, se mantiene o decrece.\par
\campo{Condiciones:} $p\in\{4,5,6,7\}$, $q\in\{1,2,3\}$ y
$p-q-1\in\{1,2,3,4\}$.\par
\campo{Control:}
\[
 \frac{c_{n+1}}{c_n}=\frac{p}{n+q+1}.
\]
Hay igualdad cuando $n=p-q-1$; evitar esa igualdad si se desea solo crecimiento
y decrecimiento eligiendo $p$ no entero.

\familia{4}{Monotonía y acotación de una raíz}
\campo{Plantilla:}
\[
 d_n=\sqrt{n+s}-\sqrt{n+r},\qquad s>r.
\]
Pedir racionalizar, probar monotonía, determinar cotas y anticipar el límite.\par
\campo{Condiciones:} $r\geq0$, $1\leq s-r\leq6$ y $r,s\leq8$.\par
\campo{Tuplas $(r,s)$:}
\[
 (0,3),\ (1,5),\ (2,6),\ (4,7).
\]
\campo{Control:}
\[
 d_n=\frac{s-r}{\sqrt{n+s}+\sqrt{n+r}},
\]
por lo que es positiva, decreciente y converge a $0$.

\familia{5}{Límite con potencia dominante oculta}
\campo{Plantilla:}
\[
 e_n=\ln\!\left(A\lambda^n+B\mu^n\right)-n\ln\lambda,
\qquad \lambda>\mu>0.
\]
Pedir factorizar la potencia dominante y calcular el límite.\par
\campo{Condiciones:} $A,B\in\{1,2,3,4\}$, $(\lambda,\mu)$ en
$\{(3,2),(4,2),(4,3),(5,2)\}$ y argumentos positivos.\par
\campo{Control:}
\[
 e_n=\ln\!\left(A+B(\mu/\lambda)^n\right)
\longrightarrow\ln A.
\]
Para evitar respuesta nula, escoger $A\neq1$.

\familia{6}{Iteración de Herón}
\campo{Plantilla:}
\[
 x_1=u,\qquad
 x_{n+1}=\frac12\left(x_n+\frac{Q}{x_n}\right).
\]
Pedir probar positividad, acotación, monotonía y convergencia.\par
\campo{Condiciones:} $Q$ no cuadrado perfecto, $Q\in\{2,3,5,6,7\}$ y
$u$ entero con $u>\sqrt Q$, procurando $u^2-Q\leq6$.\par
\campo{Tuplas $(Q,u)$:}
\[
 (2,2),\ (3,2),\ (5,3),\ (6,3),\ (7,3).
\]
\campo{Control:} para $n\geq2$, $x_n\geq\sqrt Q$; la sucesión decrece y
$\lim x_n=\sqrt Q$.

\familia{7}{Convergencia o divergencia mediante subsucesiones}
\campo{Plantilla:}
\[
 y_n=A+B(-1)^n+\frac{C}{n}.
\]
Pedir estudiar las subsucesiones pares e impares y decidir convergencia.\par
\campo{Condiciones:} $A,C\in\mathbb Z$, $B\in\{-3,-2,-1,1,2,3\}$,
$|A|,|C|\leq5$. Mantener $B\neq0$ para garantizar dos límites distintos.\par
\campo{Tuplas $(A,B,C)$:}
\[
 (1,1,2),\ (2,-1,3),\ (-1,2,1),\ (0,-2,5).
\]
\campo{Control:}
\[
 y_{2n}\longrightarrow A+B,
\qquad
 y_{2n-1}\longrightarrow A-B;
\]
por tanto, $(y_n)$ diverge.

\end{document}
